\documentclass[a4paper,11pt]{article}
\usepackage[utf8]{inputenc}
\usepackage[T1]{fontenc}
\usepackage[french]{babel}
\usepackage{lmodern}
\usepackage{amsmath,amssymb,amsthm}
\usepackage{mathrsfs}
\usepackage{enumitem}
\usepackage{multicol}

\setlist[enumerate]{label=\textbf{\textcolor{Couleur8}{\arabic*.}}, left=1em}

% Si vous voulez aussi modifier les listes enumerate imbriquées :
\setlist[enumerate,2]{label=\textcolor{Couleur8}{\alph*)}}
\setlist[enumerate,3]{label=\textcolor{Couleur8}{\roman*)}}
\setlist[enumerate,4]{label=\textcolor{Couleur8}{\Alph*)}}
\usepackage{graphicx}
\usepackage{xcolor}
\usepackage{tcolorbox}
\usepackage{geometry}
\usepackage{fancyhdr}
\usepackage{titlesec}
\usepackage{cancel}
\usepackage{ifthen}
\usepackage{tikz}
\usetikzlibrary{arrows.meta}
\usepackage{tkz-tab}
\usepackage{eurosym}

% OPTION PROF/ÉLÈVE
% Décommentez UNE des deux lignes suivantes :
\newboolean{prof}\setboolean{prof}{true}   % Version professeur (avec corrections des devoirs)
\newboolean{prof}\setboolean{prof}{true}  % Version élève (sans corrections des devoirs)

\definecolor{Couleur1}{HTML}{4A5D7A} %Th
\definecolor{Couleur2}{HTML}{A5B5D6} %Prop
\definecolor{Couleur3}{HTML}{A8C68A} %Méthode
\definecolor{Couleur6}{HTML}{8A9CC1} %Def
\definecolor{Couleur7}{HTML}{E6B459} %Rq Voc
\definecolor{Couleur8}{HTML}{D36740} %Titres
\definecolor{correction}{HTML}{D36740} % Couleur pour les corrections
\definecolor{coopmathscorr}{HTML}{F15929} % Couleur corrections coopmaths

% Configuration des marges
\geometry{
	top=2cm,
	bottom=2cm,
	inner=1.5cm,
	outer=1.5cm,
}

% Police
\renewcommand{\familydefault}{\sfdefault}

% Compteurs
\newcounter{seancenum}
\newcounter{seancenumD}
\setcounter{seancenum}{1}
\setcounter{seancenumD}{1}

% Configuration des en-têtes et pieds de page
\pagestyle{fancy}
\fancyhf{}
\ifthenelse{\boolean{prof}}{
    \fancyhead[L]{\textcolor{Couleur8}{\textbf{Corrections Automatismes - Classe de seconde - Version Professeur}}}
}{
    \fancyhead[L]{\textcolor{Couleur8}{\textbf{Corrections Devoirs - Séance 4 - Classe de seconde}}}
}
\fancyhead[R]{\textcolor{Couleur8}{\textbf{2025-2026}}}
\fancyfoot[L]{\textcolor{Couleur8}{Mme Bertail et Mme Pinto}}
\fancyfoot[R]{\textcolor{Couleur8}{\thepage}}
\renewcommand{\headrulewidth}{1pt}
\renewcommand{\headrule}{\hbox to\headwidth{\color{Couleur8}\leaders\hrule height \headrulewidth\hfill}}

% Environnements personnalisés
\newtcolorbox{seance}[1]{
    colback=Couleur6!10,
    colframe=Couleur1!65,
    fonttitle=\bfseries\large,
    title=Correction Séance \theseancenum #1,
    left=5mm,
    right=5mm,
    top=3mm,
    bottom=3mm,
    arc=0mm,
    after=\stepcounter{seancenum}
}

\newtcolorbox{seanceD}[1]{
    colback=Couleur6!5,
    colframe=Couleur6!45,
    fonttitle=\bfseries\large,
    title=Correction Devoirs - Séance \theseancenumD #1,
    left=5mm,
    right=5mm,
    top=3mm,
    bottom=3mm,
    arc=0mm,
    after=\stepcounter{seancenumD}
}

\begin{document}

\setcounter{seancenumD}{4}
\newpage
\begin{seanceD}{} %4

\begin{enumerate}
\item La seule égalité vraie est : ${\color{correction}\boldsymbol{7^{-4}\times 4^{-4}=28^{-4}}}$.
En effet, $7^{-4}\times 4^{-4}=(7\times 4)^{-4}=28^{-4}$

Concernant les autres propositions :
$\dfrac{3^{-5}}{3^{12}}=3^{-17}\neq 3^{7}$
$20\times \dfrac{1}{20^{3}}=20^{-2}\neq 20^{2}$
$\left(20^{-4}\right)^4=20^{-16}\neq 20^{0}$

\item $N=\dfrac{10^5}{2^2}=\dfrac{2^5\times 5^5 }{2^2}=5^5\times 2^{3}=10^3\times 5^{2}={\color{correction}\boldsymbol{25\times 10^{3}}}$

\item On utilise la formule $a^n\times a^m=a^{n+m}$ avec $a=-6$, $n=-2$ et $p=-9$.
$(-6)^{-2}\times (-6)^{-9}=(-6)^{-2+(-9)}={\color{correction}\boldsymbol{(-6)^{-11}}}$

\item On applique la propriété des puissances de puissances d'un réel : \\
    Soit $n\in \mathbb{N}$, et $p \in \mathbb{N}$, on a :  \\
     $\left(a^{n}\right)^{p}=a^{np}$\\
    $\begin{aligned}\left(4^{n}\right)^{3}&=\left(4\right)^{3n}\\
    &=\left(4^{3}\right)^{n}\\
    &=64^{n}
    \end{aligned}$\\La bonne réponse est la réponse {\bfseries \color{correction}C}.

\item De la relation $7z+8c=7$, on déduit en ajoutant $-8c$ dans chaque membre :
$7z=7-8c$
Puis, en divisant par $7$, on obtient : ${\color{correction}\boldsymbol{z=\dfrac{7-8c}{7}}}$.


\item La réunion $]-21\,;\,14]\,\cup \,[ -28\,;\,-6]$ = ${\color{correction}\boldsymbol{[-28\,;\,14]}}$ car les deux intervalles se chevauchent.


\end{enumerate}

\end{seanceD}

\end{document}